% from ega0-19-fr.tex
\oldpage[0\textsubscript{IV-1}]{97}

\begin{lemma}[19.5.6.5]
\label{0.19.5.6.5-fr}
Supposons que les $A$-algèbres $B$ et $C$ soient formellement lisses et que
$B$ soit préadmissible (de sorte qu'en vertu de (19.5.6.1), les conditions de
(19.5.6.3) sont satisfaites). Pour tout système de deux indices
$\alpha\leq\beta$ et d'un homomorphisme $v_{\alpha\beta}$ vérifiant les
conditions de (19.5.6.3, (i)), il existe un indice $\lambda\geq\beta$ et un
$A$-homomorphisme d'algèbres
\[
  w_{\beta\lambda}:E_{\lambda,n}\longrightarrow F_{\beta,n}
\]
tel que~: $1^\circ$ $\gr^0(w_{\beta\lambda}):C_\lambda\to C_\beta$ et
$\gr^1(w_{\beta\lambda}):
\mathfrak J_\lambda/\mathfrak J_\lambda^2\to
\mathfrak J_\beta/\mathfrak J_\beta^2$ soient les homomorphismes de
transition~; $2^\circ$ le composé
$E_{\lambda,n}\xrightarrow{w_{\beta\lambda}}F_{\beta,n}
\xrightarrow{v_{\alpha\beta}}E_{\alpha,n}$ soit l'homomorphisme de transition.
\end{lemma}

\begin{proof}
Appliquons le lemme (19.5.6.3, (ii)) avec $\gamma=\beta$, ce qui donne un
$\delta\geq\beta$ et un
$v_{\beta\delta}:F_{\delta,n}\to E_{\beta,n}$. Rappelons que
$v_{\beta\delta}$ est surjectif~; d'autre part, son noyau $\mathfrak R$ est
\emph{nilpotent}~: en effet, l'idéal d'augmentation $F_{\delta,n}^+$ de
$F_{\delta,n}$ est nilpotent par définition, et le noyau de
$\gr^0(v_{\beta\delta}):C_\delta\to C_\beta$ est aussi nilpotent en vertu de
l'hypothèse que $B$ est préadmissible. On voit donc que $E_{\beta,n}$
s'identifie à $F_{\delta,n}/\mathfrak R$. Comme $B$ est une $A$-algèbre
formellement lisse, le $A$-homomorphisme canonique
$p_{\beta,n}:B\to E_{\beta,n}$, qui est continu, se factorise en
$B\xrightarrow{w}F_{\delta,n}\xrightarrow{v_{\beta\delta}}E_{\beta,n}$, où
$w$ est un $A$-homomorphisme continu. Comme $F_{\delta,n}$ est discret, il
existe $\lambda\geq\delta$ tel que $w$ soit nul dans $\mathfrak b_\lambda$,
donc $w$ se factorise en
$B\to B/\mathfrak b_\lambda\xrightarrow{w'_\lambda}F_{\delta,n}$. Considérons
l'homomorphisme composé
$B/\mathfrak b_\lambda\xrightarrow{w'_\lambda}F_{\delta,n}
\xrightarrow{q_{\beta\delta}}F_{\beta,n}$, où $q_{\beta\delta}$ est
l'homomorphisme de transition~; notons que
$\gr^0(q_{\beta\delta})=\gr^0(v_{\beta\delta}):C_\delta\to C_\beta$ est
l'homomorphisme de transition~; comme le composé
$v_{\alpha\beta}\circ q_{\beta\delta}\circ w'_\lambda$ est l'homomorphisme
canonique $B/\mathfrak b_\lambda\to
B/(\mathfrak b_\alpha+\mathfrak J^{n+1})$, cela montre que l'image de
$\mathfrak J_\lambda=(\mathfrak b_\lambda+\mathfrak J)/\mathfrak b_\lambda$
par $q_{\beta\delta}\circ w'_\lambda$ est contenue dans l'idéal d'augmentation
$F_{\beta,n}^+$, et par suite l'image de
$\mathfrak J_\lambda^{n+1}=(\mathfrak b_\lambda+\mathfrak J^{n+1})/
\mathfrak b_\lambda$ est nulle. Autrement dit,
$q_{\beta\delta}\circ w'_\lambda$ se factorise en
\[
  B/\mathfrak b_\lambda\to
  B/(\mathfrak b_\lambda+\mathfrak J^{n+1})=E_{\lambda,n}
  \xrightarrow{w_{\beta\lambda}}F_{\beta,n}
\]
tel que le composé
$E_{\lambda,n}\xrightarrow{w_{\beta\lambda}}F_{\beta,n}
\xrightarrow{v_{\alpha\beta}}E_{\alpha,n}$ soit l'homomorphisme de transition.
Le raisonnement précédent montre aussi que $\gr^0(w_{\beta\lambda})$, qui est
le composé
$B/(\mathfrak b_\lambda+\mathfrak J)\to\gr^0(F_{\delta,n})
\xrightarrow{\gr^0(q_{\beta\delta})}\gr^0(F_{\beta,n})$, est l'homomorphisme
de transition $C_\lambda\to C_\beta$, puisque
$v_{\beta\delta}\circ w$ est l'homomorphisme canonique. En outre, on a aussi
$\gr^1(q_{\beta\delta})=\gr^1(v_{\beta\delta})$ (19.5.6.3), donc le même
raisonnement prouve que $\gr^1(w_{\beta\lambda})$ est l'homomorphisme de
transition.
\end{proof}

\begin{env}[19.5.6.6]
\label{0.19.5.6.6-fr}
\emph{Démonstration de} (19.5.3, (ii)). --- Pour montrer que $\varphi_n$ est un
bimorphisme formel, nous allons appliquer le critère de (19.1.3, (iii)). Les
conditions du lemme (19.5.6.5) étant satisfaites par hypothèse, déterminons,
pour tout indice $\alpha$, $v_{\alpha\beta}$ et $w_{\beta\lambda}$ vérifiant
les conclusions de ce lemme. L'homomorphisme
\[
  \gr^n(v_{\alpha\beta}):\gr^n(F_{\beta,n})
  \longrightarrow\gr^n(E_{\alpha,n})
\]
n'est autre que l'homomorphisme
\[
  \varphi_{\alpha\beta,n}:
  \mathbf S_{C_\beta}^n(\mathfrak J_\beta/\mathfrak J_\beta^2)
  \longrightarrow\mathfrak J_\alpha^n/\mathfrak J_\alpha^{n+1}
\]
déduit de l'homomorphisme canonique $\varphi_n$ (19.5.2.1) par passage aux
quotients~; en effet, il résulte de (19.5.6.3) que $\gr^0(v_{\alpha\beta})$ et
$\gr^1(v_{\alpha\beta})$ coïncident respectivement avec
$\varphi_{\alpha\beta,0}$ et $\varphi_{\alpha\beta,1}$, et la définition de
l'homomorphisme canonique $\varphi$ montre alors, par récurrence sur $j\leq n$,
que $\gr^j(v_{\alpha\beta})$ et $\varphi_{\alpha\beta,j}$ sont égaux pour tout
$j$. Cela étant, comme $\varphi_n$ est surjectif, c'est \emph{a fortiori} un
épimorphisme formel~; en outre, pour $j\leq n$, l'homomorphisme composé
\[
  \gr^j(F_{\lambda,n})\xrightarrow{\varphi_{\lambda\lambda,j}}
  \gr^j(E_{\lambda,n})\xrightarrow{\gr^j(w_{\beta\lambda})}
  \gr^j(F_{\beta,n})
\]
est l'homomorphisme de transition, car c'est vrai pour $j=0$ et $j=1$ en
vertu de (19.5.6.5), et comme $\gr^*(E_{\lambda,n})$ est engendré par
$\gr^1(E_{\lambda,n})$ cela démontre l'assertion par récurrence sur $j$.
Composant avec l'homomorphisme de transition
$\gr^j(F_{\beta,n})\to\gr^j(F_{\alpha,n})$, on voit donc, pour $j=n$, que
l'on a factorisé l'homomorphisme de transition
$\gr^n(F_{\lambda,n})\to\gr^n(F_{\alpha,n})$ en
\[
  \gr^n(F_{\lambda,n})\xrightarrow{\varphi_{\lambda\lambda,n}}
  \gr^n(E_{\lambda,n})\longrightarrow\gr^n(F_{\alpha,n}),
\]
ce qui est la condition du critère (19.1.3) pour que $\varphi_n$ soit un
bimorphisme formel.
\end{env}

